% Báo cáo kỹ thuật đồ án Project 3 — Truy vấn đa phương thức có giải thích bằng đồ thị (nhóm 7).
% Biên dịch: bash build.sh (XeLaTeX + BibTeX, chạy 3 lượt). Mỗi mục nằm trong sections/; mỗi con số trong báo cáo
% lấy từ một file kết quả trong experiments/ (đường dẫn ghi ngay tại chỗ hoặc ở cuối mục).
\documentclass[conference,a4paper]{IEEEtran}

\usepackage{amsmath,amssymb}
\usepackage{newtxmath} % công thức kiểu Times cho khớp chữ
\usepackage[no-math]{fontspec}
% TeX Gyre Termes (bản sao của Times, có chữ in hoa nhỏ đủ dấu tiếng Việt) để mọi máy build ra giống nhau.
\IfFontExistsTF{texgyretermes-regular.otf}{\setmainfont{texgyretermes}[Extension=.otf, UprightFont=*-regular,
   BoldFont=*-bold, ItalicFont=*-italic, BoldItalicFont=*-bolditalic]}{\setmainfont{Times New Roman}}
\usepackage{polyglossia}
\setdefaultlanguage{vietnamese}
\usepackage{graphicx}
\usepackage{booktabs}
\usepackage{array}
\usepackage{multirow}
\usepackage{xcolor}
\usepackage{tikz}
\usetikzlibrary{arrows.meta,positioning,calc,fit,backgrounds}
\usepackage{pgfplots}
\pgfplotsset{compat=1.18, /pgf/number format/use comma, /pgf/number format/1000 sep={}}
\usepackage{enumitem}
\setlist{nosep,leftmargin=*}
\usepackage[hyphens]{url}
\usepackage[colorlinks, urlcolor=blue, linkcolor=black, citecolor=black]{hyperref}
\usepackage{etoolbox}
\gappto\captionsvietnamese{\renewcommand{\refname}{Tài liệu tham khảo}}
\renewcommand{\IEEEkeywordsname}{Từ khoá}
\emergencystretch=2em

\graphicspath{{figures/}}

% Màu dùng chung cho hình tự vẽ
\definecolor{cClip}{HTML}{2A6F97}   % kênh CLIP
\definecolor{cObj}{HTML}{2A9D8F}    % kênh vật thể / thực thể
\definecolor{cTri}{HTML}{C8553D}    % kênh bộ ba / cạnh
\definecolor{cGray}{HTML}{5C6570}

% Chỗ còn thiếu số hoặc cần nhóm xác nhận: in đỏ để không lọt vào bản nộp.
\newcommand{\todo}[1]{\textcolor{red}{[#1]}}
% Số thập phân trong chế độ toán viết 0{,}6 để không có khoảng trắng sau dấu phẩy.
\newcommand{\pmstd}[2]{#1\,$\pm$\,#2}  % mean ± std trong bảng (chế độ văn bản)
\newcommand{\file}[1]{{\small\nolinkurl{#1}}} % đường dẫn file trong repo; gõ nguyên dấu _, không escape
% Nguồn số liệu của một bảng, in ngay dưới bảng (caption của IEEEtran in hoa nên không đặt đường dẫn ở đó).
\newcommand{\tabsrc}[1]{\par\smallskip{\scriptsize\raggedright Nguồn: #1\par}}
\newcommand{\repo}{\url{https://github.com/HatDuaa/multimodal-graph-retrieval-xai}}

\begin{document}

\title{Truy vấn văn bản--ảnh có giải thích bằng đồ thị:\\
scene graph, đồ thị tri thức đa phương thức và đánh giá lời giải thích}

\author{
\IEEEauthorblockN{Nhóm 7: Nguyễn Đình Lộc (25C11050), Ngô Trương Minh Đạt (25C15003), Trần Đắc Khoa (25C15015),\\
Nguyễn Bùi Mẫn Nhi (25C15055), Âu Dương Khang (25C15045)}
\IEEEauthorblockA{Khoa Công nghệ Thông tin, Trường Đại học Khoa học Tự nhiên, ĐHQG-HCM\\
Học phần: Học máy với dữ liệu đồ thị \quad Giảng viên: TS. Lê Ngọc Thành\\
Mã nguồn: \repo}
}

\maketitle

\begin{abstract}
Báo cáo trình bày một hệ thống truy vấn văn bản--ảnh (text-to-image retrieval). Hệ thống dùng dữ liệu đồ thị để xếp
hạng lại kết quả của CLIP đóng băng, và kèm mỗi kết quả một lời giải thích lấy từ chính phép tính xếp hạng. Trên Visual
Genome $\cap$ COCO (cấp độ~1), một mô hình GATv2 chạy trên scene graph của ảnh và trên đồ thị tách từ câu truy vấn. Mô
hình nâng R@1 trên tập test từ 39,32 lên \pmstd{47,95}{0,15} (3 seed). Huấn luyện lại sau khi bỏ mọi cạnh làm R@1 val giảm
2,39 điểm; thay đồ thị của ảnh bằng đồ thị của một ảnh khác làm R@1 val giảm xuống 21,91, thấp hơn CLIP thuần. Trên MKG-W (cấp độ~2),
nhóm tái lập NativE (SIGIR 2024), đạt MRR 36,99 so với 36,58 của bài báo, rồi đưa đồ thị tri thức vào điểm xếp hạng. Kênh
đếm cạnh của đồ thị tri thức nâng R@1 test từ 24,89 lên 29,82. Trên 15 truy vấn gián tiếp tự viết, R@1 tăng từ 40 lên 80.
Khi kết hợp với kênh đếm cạnh, tập validation chọn trọng số 0 cho tín hiệu của NativE. Ở cấp độ~3, nhóm đánh giá lời giải
thích của mô hình chính trên toàn bộ tập validation bằng cách bỏ hoặc chỉ giữ một số phần của đồ thị rồi chấm lại: bỏ 3 phần
đứng đầu lời giải thích làm ảnh mất hạng 1 ở 58,1\% số câu, so với 9,4\% khi bỏ 3 phần ngẫu nhiên và 34,9\% với mô hình
cấp~1 gốc trộn điểm qua z-score. Trong các cách giải thích không ước lượng giá trị Shapley, lời giải thích này khớp Shapley
tốt nhất mà không cần chấm thêm lần nào. Nhóm cũng tự xây một đồ thị tri thức đa phương thức tiếng Việt
(150 thực thể, 325 triple, 331 ảnh); chạy nguyên phương pháp của cấp độ~2 trên đó, không chọn lại tham số nào, nâng R@1 từ
23,35 lên 33,92. Ở mọi thực nghiệm, trọng số kết hợp, checkpoint và từng bước xây mô hình được chọn trên train hoặc
validation, không dựa trên tập test. Kết quả báo cáo dạng trung bình $\pm$ độ lệch chuẩn qua 3 seed.
\end{abstract}

\begin{IEEEkeywords}
truy vấn đa phương thức, scene graph, đồ thị tri thức đa phương thức, mạng nơ-ron đồ thị, xếp hạng lại, giải thích mô
hình, fidelity
\end{IEEEkeywords}

% =====================================================================
\section{Giới thiệu}
\label{sec:intro}

CLIP~\cite{radford2021clip} đưa ảnh và câu vào cùng một không gian vector, nên truy vấn văn bản--ảnh chỉ còn là tìm ảnh
có cosine lớn nhất với câu. Cách so khớp này có hai giới hạn. Thứ nhất, nó không dùng quan hệ có cấu trúc giữa các vật
thể trong ảnh, hay giữa các thực thể trong một đồ thị tri thức. Thứ hai, nó không cho biết vì sao một ảnh được xếp lên
đầu. Đề bài yêu cầu một hệ thống trả lời ba câu hỏi nối tiếp nhau:
(1) ảnh nào phù hợp nhất với truy vấn;
(2) tri thức quan hệ giúp truy vấn tốt hơn ở đâu, nhất là với truy vấn gián tiếp;
(3) vì sao mô hình trả về kết quả này.
Nhóm làm nhánh A của đề (văn bản tìm ảnh) và thực hiện đủ ba cấp độ.

\textbf{Cấp độ~1} (Visual Genome $\cap$ COCO). CLIP đóng băng chọn 50 ứng viên. Một bộ mã hoá GATv2~\cite{brody2022gatv2}
chạy trong scene graph của từng ảnh và trong đồ thị tách từ câu, rồi chấm điểm khớp ở hai mức: vật thể và bộ ba (chủ thể,
quan hệ, đối tượng). Điểm cuối kết hợp cosine CLIP với điểm đồ thị theo dạng
$\alpha\cdot\text{CLIP} + (1-\alpha)\cdot\text{đồ thị}$ của đề. Nhóm so với CLIP thuần trên cùng cách chia. Các ablation
gồm bỏ cạnh, nối lại cạnh ngẫu nhiên và thay đồ thị của ảnh khác. Phần cuối phân tích 5 truy vấn thành công và 5 truy vấn
thất bại.

\textbf{Cấp độ~2} (thêm MKG-W). Nhóm tái lập NativE~\cite{zhang2024native}, một mô hình hoàn thiện đồ thị tri thức đa
phương thức năm 2024, và kiểm tra lại con số link prediction của bài báo. Sau đó nhóm đưa ba loại tín hiệu đồ thị vào cùng
pipeline xếp hạng lại: số cạnh tri thức, embedding của NativE, và điểm bộ ba của chính NativE. Nhóm xây thêm bộ 25 truy vấn
tự viết, tách nhóm trực tiếp và gián tiếp.

\textbf{Cấp độ~3}. Nhóm đo fidelity, validity và sparsity của lời giải thích trên toàn bộ tập validation, kèm bằng chứng
trên từng truy vấn. Nhóm cũng xây một đồ thị tri thức đa phương thức (MMKG) tiếng Việt nhỏ, chạy phương pháp của cấp độ~2
trên đó và phân tích khả năng chuyển giao.

Mục~\ref{sec:data} mô tả dữ liệu và các quy tắc chống rò rỉ. Mục~\ref{sec:method} mô tả mô hình và cách lời giải thích được sinh ra. Các
mục~\ref{sec:l1}--\ref{sec:vnmmkg} trình bày kết quả theo từng cấp độ. Mục~\ref{sec:conclusion} nêu hạn chế. Mục~\ref{sec:contrib}
ghi phân công và phần đã dùng công cụ AI. Mọi con số đều lấy từ một file kết quả trong thư mục \file{experiments/} của
repo; đường dẫn ghi ngay cạnh bảng.

% =====================================================================
\section{Dữ liệu, cách chia tập và chống rò rỉ}
\label{sec:data}

\subsection{Visual Genome $\cap$ COCO (cấp độ 1)}
\label{sec:data-vg}

Tập công khai thứ nhất là phần giao giữa Visual Genome~\cite{krishna2017vg} (VG, ảnh có scene graph gán tay: vật thể,
thuộc tính, quan hệ) và MS-COCO~\cite{lin2014coco} (mỗi ảnh có 5 caption). Hai bộ nối với nhau theo \texttt{coco\_id}. Truy
vấn là caption COCO; đáp án là ảnh chứa caption đó. Nhóm chia theo split Karpathy~\cite{karpathy2015deep}, cách chia chuẩn
của bài toán truy vấn ảnh--văn bản trên COCO, và gộp phần \texttt{restval} vào train. Nhóm lấy toàn bộ phần giao, không
lấy mẫu, và bỏ 290 ảnh VG trùng \texttt{coco\_id} với ảnh khác (Bảng~\ref{tab:data}). Khi đánh giá một split, tập ứng viên
(pool) là toàn bộ ảnh của chính split đó. Pool test (2\,138 ảnh) nhỏ hơn MSCOCO 5K, nên số của nhóm không so trực tiếp được
với số trong các bài báo.

\textbf{Đồ thị ảnh.} Mỗi ảnh có một scene graph riêng, lấy từ \texttt{objects.json} và \texttt{relationships.json} của VG.
Nhãn giữ dạng mở, chỉ chuẩn hoá alias, không lọc về VG150 (bước kiểm tra độ phủ cho thấy lọc VG150 làm mất các nhãn hay
gặp như \emph{water}, \emph{grass}, \emph{field}). Quan hệ có hướng; các quan hệ trùng cả ba thành phần được gộp. Không có
cạnh giữa các ảnh, nên đồ thị không chứa cạnh suy ra từ thống kê toàn tập. Mỗi ảnh train và val có trung bình 26,7 node và
18,0 cạnh; 3,0\% ảnh không có cạnh nào.

\textbf{Đồ thị câu.} Mô hình cấp độ~1 tách caption bằng SceneGraphParser~\cite{schuster2015sgparser} (gói
\texttt{sng\_parser}): trung bình 3,06 vật thể và 1,65 quan hệ mỗi câu, và 14,2\% câu val không tách được quan hệ nào. Mô
hình chính dùng FACTUAL-T5 large~\cite{li2023factual}, một parser seq2seq đã huấn luyện sẵn để sinh scene graph từ câu.
Parser chỉ chạy suy luận, không tinh chỉnh. Trên val, FACTUAL giảm tỉ lệ node là đại từ từ 3,9\% xuống 0,1\%, và tăng tỉ lệ
câu có ít nhất một quan hệ từ 85,8\% lên 97,3\%. (Nguồn: \file{experiments/data_stats.md},
\file{experiments/improve/parser-compare/REPORT.md}.)

\begin{table}[t]
\centering
\caption{Kích thước hai tập công khai. Mỗi thực thể MKG-W có đúng một ảnh và một truy vấn.}
\label{tab:data}
\small
\begin{tabular}{lrrr}
\toprule
 & Train & Val & Test \\
\midrule
VG $\cap$ COCO: ảnh & 46\,944 & 2\,126 & 2\,138 \\
VG $\cap$ COCO: caption (truy vấn) & 234\,854 & 10\,633 & 10\,696 \\
\midrule
MKG-W retrieval: thực thể & 7\,136 & 892 & 892 \\
MKG-W KGC: triple & 34\,196 & 4\,276 & 4\,274 \\
\bottomrule
\end{tabular}
\end{table}

\subsection{MKG-W (cấp độ 2)}
\label{sec:data-mkgw}

Tập công khai thứ hai là MKG-W~\cite{xu2022mmrns}, một đồ thị tri thức đa phương thức lấy từ Wikidata~\cite{vrandecic2014wikidata}:
mỗi thực thể có nhãn, mô tả và ảnh. Nhóm dùng nó cho hai việc trên cùng một bộ thực thể.
\begin{itemize}
\item \textbf{Đồ thị tri thức.} Bản MKG-W của NativE có 15\,000 thực thể, 169 quan hệ, và các triple chia sẵn cho bài toán
link prediction (KGC) như ở Bảng~\ref{tab:data}.
\item \textbf{Bài toán truy vấn.} Chỉ xét thực thể có cả mô tả lẫn ảnh. Truy vấn là mô tả tiếng Anh của thực thể, đã che
tên, ví dụ ``a commune in Pyrénées-Atlantiques, France''. Đáp án là ảnh của chính thực thể đó. Ảnh lấy từ bản phát hành của
nhóm tác giả MMRNS: 37\,192 ảnh của 8\,961 thực thể; mỗi thực thể dùng một ảnh. Đây cũng là bộ ảnh mà NativE dùng để trích
đặc trưng ảnh. Split retrieval cố định với seed 0; tổng 8\,920 thực thể.
\end{itemize}
Trên val, 78,81\% truy vấn liên kết được với ít nhất một thực thể của đồ thị (test 80,94\%). Với 45,6\% truy vấn val, đáp án
nối trực tiếp với một thực thể được nhắc trong câu bằng một cạnh KGC-train.

\subsection{Chống rò rỉ dữ liệu}
\label{sec:leakage}

Mục~6 của đề là tiêu chí chấm. Bảng~\ref{tab:leakage} liệt kê từng quy tắc và cách nhóm bảo đảm nó trong code hoặc quy
trình.

\begin{table*}[t]
\centering
\caption{Các quy tắc chống rò rỉ và cách bảo đảm.}
\label{tab:leakage}
\small
\begin{tabular}{p{0.30\textwidth}p{0.64\textwidth}}
\toprule
Quy tắc & Cách bảo đảm \\
\midrule
Mọi mô hình so sánh dùng cùng một tệp chia &
File split trong \file{data/splits/} là nguồn duy nhất; SHA-256 ghi ở \file{data/MANIFEST.md}. Mọi dòng trong một bảng dùng
chung danh sách top-50 của CLIP. \\
Test không dùng để huấn luyện, chọn siêu tham số hay chọn mô hình &
VG: mỗi mô hình đã chốt chấm test một lần, 25/09 (mô hình cấp~1) và 25/10 (mô hình chính); scene ranker
(Mục~\ref{sec:contrib}) cũng chấm một lần. \texttt{GraphStore} từ chối nạp test nếu không gọi rõ
\texttt{allow\_test=True}; trước khi chấm, script kiểm tra ứng viên test tái tạo đúng số của CLIP thuần.
MKG-W: lần chấm báo cáo là 17/11. Trước đó, một lần chạy thử code với embedding giả trên test đã in ra số của dòng đếm
cạnh; ngày 18/11, một thành viên chạy lại độc lập toàn bộ cấp~2 với NativE tự huấn luyện (Mục~\ref{sec:l2}). Cả hai lần
không dẫn tới thay đổi nào: $\alpha,\beta$ chỉ lấy từ val. \\
Cạnh thống kê chỉ xây từ train &
Đồ thị ảnh VG là chú thích của từng ảnh, không có cạnh giữa các ảnh. Kênh đồ thị của MKG-W chỉ nạp triple KGC-train; hàm
dựng đồ thị có \texttt{assert kgc\_split == "train"}. \\
$\alpha$ và mọi lựa chọn dựa trên train/val &
VG: trọng số kết hợp $a,b,c$ học trên train; dừng sớm, chọn checkpoint, ablation và từng bước xây mô hình chính dựa trên
val; lr và weight decay cố định từ đầu, không dò. Với MKG-W, $\alpha,\beta$ chọn trên val
cho từng seed, ghi vào \file{experiments/level2/val/choice.json}; lần chạy test chỉ đọc lại file này. \\
Ứng viên lúc huấn luyện không chạm val/test &
Train được chia cố định (seed 0) thành 22 nhóm khoảng 2\,134 ảnh. Top-50 của mỗi caption train chỉ lấy trong nhóm của ảnh
đúng, nên độ khó giống lúc đánh giá. CLIP và FACTUAL-T5 đóng băng. \\
Mô hình đồ thị tri thức không thấy truy vấn và ảnh của val/test &
Đặc trưng văn bản của NativE mã hoá chính mô tả thực thể, tức chính câu truy vấn; đặc trưng ảnh mã hoá chính ảnh ứng viên.
Bản \texttt{MKG-W-rmask} đặt về 0 hai đặc trưng này cho 1\,784 thực thể val/test retrieval trước khi huấn luyện NativE. \\
Bộ truy vấn tự viết chỉ để đánh giá &
Bộ 25 câu MKG-W được commit một lần (16/11) cùng script kiểm tra, không sửa sau đó; lần chấm dùng lại $\alpha,\beta$
của val. \\
\bottomrule
\end{tabular}
\end{table*}

\subsection{Kiểm tra pipeline}
\label{sec:checks}

Trước khi tin các con số, nhóm kiểm tra pipeline bằng những phép so có đáp án biết trước.
\begin{itemize}
\item \textbf{CLIP trên MSCOCO 5K.} Chạy chính code của nhóm (mã hoá, FAISS, module đánh giá) trên pool chuẩn 5\,000 ảnh.
Kết quả lệch số đã công bố của open\_clip~\cite{cherti2023openclip} không quá 0,03 điểm (text$\to$ảnh R@1 30,45 so với
30,44). Phép so này cũng phát hiện một lỗi: trọng số OpenAI phải nạp bằng cấu hình \texttt{ViT-B-32-quickgelu}. Nạp sai
cấu hình làm R@1 thấp hơn khoảng 1 điểm. Toàn bộ đặc trưng được trích lại sau khi sửa
(\file{experiments/checks/README.md}).
\item \textbf{Mô hình chưa huấn luyện.} Ở bước khởi tạo, mô hình cấp độ~1 phải cho đúng điểm của phép thử không huấn luyện
đã đo trước (R@1 val 43,50). Có script kiểm tra tự động cho việc này.
\item \textbf{Demo.} Gõ 200 caption val vào giao diện. Demo cho cùng 50 ứng viên, cùng đồ thị câu và cùng top-1 với pipeline
đánh giá; điểm cuối lệch tối đa $4\cdot10^{-6}$.
\item \textbf{Đặc trưng MKG-W dựng lại.} Đặc trưng CLIP gốc của MKG-W không còn truy cập được, nên nhóm trích lại bằng đúng
script trong repo. Trên val, CLIP thuần cho lại đúng R@1 24,33 như kết quả đã commit trước đó.
\end{itemize}

% =====================================================================
\section{Phương pháp}
\label{sec:method}

\subsection{Pipeline chung}

Cả ba cấp độ dùng chung một pipeline, đúng quy trình gợi ý trong đề (Hình~\ref{fig:pipeline}).
\begin{enumerate}
\item CLIP ViT-B/32 (trọng số OpenAI, đóng băng) mã hoá câu truy vấn và lấy 50 ảnh có cosine lớn nhất (FAISS~\cite{johnson2019faiss}
ở VG và demo; phép nhân ma trận trực tiếp ở MKG-W vì pool chỉ có 892 ảnh).
\item Với từng ứng viên, một mô-đun đồ thị tính điểm đồ thị. Ở VG, đó là điểm khớp giữa đồ thị câu và scene graph của ảnh.
Ở MKG-W, đó là quan hệ giữa thực thể ứng viên và các thực thể được nhắc trong câu.
\item Điểm cuối kết hợp hai điểm theo dạng $\alpha\cdot\text{CLIP} + (1-\alpha)\cdot\text{đồ thị}$ của đề. Mô hình cấp~1
và MKG-W chuẩn hoá từng điểm bằng z-score trong 50 ứng viên của truy vấn; mô hình chính của VG cộng điểm thô
(Mục~\ref{sec:method-main}). Ứng viên được sắp lại theo điểm cuối.
\item Lời giải thích của mỗi ứng viên là chính các đại lượng đã cộng vào điểm đồ thị. Không có bước tìm đường đi nào chạy
sau khi đã có thứ hạng (Mục~\ref{sec:method-expl}).
\end{enumerate}
Vì chỉ xếp lại 50 ứng viên, giới hạn trên của mọi mô hình là recall@50 của CLIP. Trên test VG, chỉ số này là 96,66\%.

\begin{figure}[t]
\centering
\begin{tikzpicture}[font=\scriptsize, node distance=3mm and 4mm,
  box/.style={draw, rounded corners=2pt, align=center, minimum height=7mm, inner sep=2pt},
  arr/.style={-{Stealth[length=1.6mm]}, thick, cGray}]
\node[box, fill=cClip!12] (q) {Truy vấn\\(câu)};
\node[box, fill=cClip!12, right=of q] (clip) {CLIP\\đóng băng};
\node[box, fill=cClip!12, right=of clip] (top) {Top-50\\ứng viên};
\node[box, fill=cObj!15, below=6mm of q] (qg) {Đồ thị câu /\\thực thể nhắc tới};
\node[box, fill=cObj!15, right=of qg] (gs) {Điểm đồ thị\\(GATv2 / KG)};
\node[box, fill=cTri!12, right=of gs] (ig) {Scene graph ảnh /\\đồ thị tri thức};
\node[box, fill=yellow!20, below=6mm of gs] (fuse) {$\alpha\cdot\text{CLIP} + (1-\alpha)\cdot\text{đồ thị}$};
\node[box, right=of fuse] (out) {Thứ hạng\\+ giải thích};
\draw[arr] (q) -- (clip); \draw[arr] (clip) -- (top);
\draw[arr] (q) -- (qg); \draw[arr] (qg) -- (gs); \draw[arr] (ig) -- (gs);
\draw[arr] (top) -- (ig);
\draw[arr] (gs) -- (fuse); \draw[arr] (top.south) |- ($(fuse.north east)+(0,0.5mm)$);
\draw[arr] (fuse) -- (out);
\end{tikzpicture}
\caption{Pipeline chung. Phần màu xanh là CLIP đóng băng. Phần đồ thị chỉ chấm lại 50 ứng viên của CLIP. Lời giải thích
là các đại lượng của chính điểm đồ thị.}
\label{fig:pipeline}
\end{figure}

\subsection{Cấp độ 1: bộ xếp lại GATv2 trên scene graph}
\label{sec:method-l1}

\textbf{Mã hoá đồ thị.} Đồ thị câu và đồ thị ảnh dùng chung một bộ mã hoá. Vector khởi đầu của mỗi node là vector CLIP
text của nhãn (câu mẫu ``a photo of a \{nhãn\}''), chiếu từ 512 xuống 256 chiều. Mỗi cạnh mang vector của nhãn quan hệ.
Hai lớp GATv2~\cite{brody2022gatv2} (4 head) cho mỗi node gom thông tin từ hàng xóm theo trọng số attention học được. Mỗi
bộ ba (chủ thể, quan hệ, đối tượng) đi qua một MLP. Vector cuối của mỗi phần là $\mathrm{L2}(f + U g)$, với $f$ là vector
ban đầu, $g$ là đầu ra GAT và $U$ khởi tạo bằng 0.

\textbf{Hai kênh khớp.} Kênh vật thể chỉ so vật thể với vật thể; kênh bộ ba chỉ so bộ ba với bộ ba. Với phần $i$ của câu
và phần $j$ cùng loại của ảnh, $\mathit{sim}_{ij}$ là cosine giữa hai vector. Mỗi phần của câu được so với mọi phần cùng
loại của ảnh qua attention:
\begin{equation}
a_{ij} = \operatorname{softmax}_j\!\left(\mathit{sim}_{ij}/\tau\right),\qquad s_i = \textstyle\sum_j a_{ij}\,\mathit{sim}_{ij}.
\end{equation}
Trọng số của từng phần trong câu là
$w_i = \operatorname{softmax}_i(u^\top p_i)$, với $p_i$ là vector của phần $i$. Điểm của kênh là
$S = \sum_i w_i s_i$.

\textbf{Trộn điểm.} Ba trọng số kênh tính riêng cho từng câu từ vector CLIP của câu:
$[a, b, c] = \operatorname{softmax}(W q + \mathit{bias})$. Điểm đồ thị là
$G = z\big((b\,z(S_\text{obj}) + c\,z(S_\text{tri}))/(b+c)\big)$, và điểm cuối là $a\,z(\text{CLIP}) + (1-a)\,G$. Đây
đúng là công thức của đề với $\alpha = a$, chỉ khác ở chỗ $\alpha$ thay đổi theo câu. Mọi tham số mới khởi tạo sao cho mô
hình chưa huấn luyện cho đúng phép thử không huấn luyện đã đo trước ($U=0$, $u=0$, $W=0$,
$\mathit{bias} = \log(0{,}60;\,0{,}28;\,0{,}12)$, $\tau = 0{,}05$). Mô hình có 1\,054\,214 tham số huấn luyện được; CLIP
không nằm trong số này (chi tiết: \file{docs/level1-model.md}).

\textbf{Huấn luyện.} Loss là cross-entropy trên 50 ứng viên của điểm cuối, chia cho nhiệt độ học được. AdamW, lr
$10^{-3}$, weight decay $10^{-4}$, batch 256, dropout 0,1. Epoch đầu khoá $W$ và $\mathit{bias}$ để phần đồ thị học trước.
Dừng sớm khi R@1 val không tăng sau 2 epoch, tối đa 10 epoch, giữ checkpoint tốt nhất theo val. Mỗi cấu hình chạy 3 seed
(0, 1, 2).

\subsection{Mô hình chính của VG}
\label{sec:method-main}

Mô hình chính (\texttt{countonly\_factual\_r3}) giữ nguyên khung trên và thay bốn chỗ. Mỗi thay đổi được giữ lại vì làm R@1
val tăng qua 3 seed (Bảng~\ref{tab:l1-chain}).
\begin{enumerate}
\item \textbf{Đồ thị câu từ FACTUAL-T5} thay cho \texttt{sng\_parser}.
\item \textbf{Trộn điểm thô:} điểm cuối $= a\cdot\cos + b\cdot S_\text{obj} + c\cdot S_\text{tri}$, bỏ z-score. Ứng viên
không có phần nào để so ở một kênh nhận điểm trung bình của các ứng viên có phần đó. Nhờ bỏ chuẩn hoá, điểm cuối tách
\emph{chính xác} thành tổng các phần, nên lời giải thích không cần xấp xỉ.
\item \textbf{Đồ thị ảnh sạch và liền.} Hai node cùng loại (cùng synset hoặc hypernym trong WordNet) có IoU $\ge 0{,}7$ được
gộp làm một. Mỗi cạnh có hộp ở hai đầu được gán 6 số hình học (độ lệch tâm, tỉ lệ kích thước, IoU, tỉ lệ diện tích). Mỗi
node nhận thêm cạnh từ 5 node gần nhất, và một node tổng nối tới mọi node. Cạnh thêm này chỉ truyền tin trong GAT, không
được tính là bộ ba. Sau bước này 100\% đồ thị ảnh val liền một khối (cấp độ~1: 4,7\%).
\item \textbf{Số lượng.} Mỗi node ảnh biết có bao nhiêu vật cùng loại trong ảnh (nhóm 1 / 2 / 3--4 / $\ge5$). Luật viết tay
trên câu gán mỗi vật thể một lớp số lượng (số nêu rõ, ``a pair of'', danh từ số nhiều, \dots; còn lại là ``không rõ''). Một
bảng học được $g$ ($7\times4$, khởi tạo 0, hàng ``không rõ'' cố định bằng 0) được cộng vào $\mathit{sim}_{ij}$ của kênh vật
thể.
\end{enumerate}

\subsection{Cấp độ 2: NativE và điểm đồ thị tri thức}
\label{sec:method-l2}

\textbf{NativE.} NativE~\cite{zhang2024native} là một mô hình hoàn thiện đồ thị tri thức đa phương thức (MMKGC: dự đoán
cạnh còn thiếu bằng cả cấu trúc lẫn ảnh và chữ của thực thể). Lõi cấu trúc là AdvRelRotatE, một biến thể của
RotatE~\cite{sun2019rotate}: mỗi quan hệ là một phép quay trong không gian phức, và mẫu âm khó được sinh bằng huấn luyện đối
kháng. Đặc trưng ảnh và chữ được trộn vào embedding thực thể với trọng số phụ thuộc quan hệ, nên mô hình chịu được thực
thể thiếu ảnh hoặc thiếu chữ. Nhóm dùng NativE theo hai bước.
\begin{itemize}
\item \textbf{Bước a: tái lập.} Huấn luyện trên MKG-W gốc với đúng mã nguồn (commit \texttt{c72291c}) và siêu tham số của
tác giả (1\,000 epoch, dim 250, margin 4, 128 mẫu âm, lr $10^{-4}$, batch 1\,024). So MRR và Hits@1 với bài báo.
\item \textbf{Bước b: đưa vào truy vấn.} Huấn luyện trên bản \texttt{MKG-W-rmask} (Mục~\ref{sec:leakage}) với 3 seed, rồi dùng
mô hình đã học làm tín hiệu đồ thị.
\end{itemize}

\textbf{Liên kết truy vấn với đồ thị.} Câu truy vấn đã che tên đáp án nhưng vẫn hay nhắc tên thực thể khác (``a rock
guitarist who was a member of \emph{Pink Floyd}''). \texttt{EntityLinker} tìm trong câu mọi nhãn và alias của thực thể
trong đồ thị, ưu tiên cụm khớp dài nhất (tối đa 5 token). Gọi $L$ là tập thực thể khớp được, $c$ là thực thể ứng viên.

\textbf{Ba kênh điểm đồ thị.}
\begin{itemize}
\item \texttt{triple}: số cạnh KGC-train trực tiếp giữa $c$ và các thực thể trong $L$. Kênh này không cần huấn luyện.
\item \texttt{emb}: cosine lớn nhất giữa embedding cấu trúc NativE của $c$ và của các thực thể trong $L$.
\item \texttt{kgc}: điểm bộ ba của chính NativE, lấy giá trị lớn nhất qua mọi $l \in L\setminus\{c\}$, mọi quan hệ (169)
và cả hai chiều $(c, r, l)$, $(l, r, c)$. Khác \texttt{emb} ở chỗ có dùng quan hệ. Code tính lại đúng hàm
\texttt{forward()} của NativE và có test đối chiếu với checkpoint thật.
\end{itemize}
Câu không khớp thực thể nào nhận điểm đồ thị 0, nên thứ hạng của nó giữ nguyên như CLIP. Điểm cuối là
\begin{equation}
\alpha\,z(\text{CLIP}) + (1-\alpha)\,z\big(\beta\,z(N) + (1-\beta)\,z(\texttt{triple})\big),
\end{equation}
với $N$ là \texttt{emb} hoặc \texttt{kgc}. Lưới tìm $\alpha \in \{0; 0{,}05; \dots; 1\}$,
$\beta \in \{0; 0{,}25; 0{,}5; 0{,}75; 1\}$, chọn trên val riêng cho từng seed theo R@1 (hoà thì xét MRR).

\subsection{Lời giải thích sinh từ cơ chế xếp hạng}
\label{sec:method-expl}

Đề yêu cầu lời giải thích phải đến từ chính cơ chế xếp hạng. Nhóm không chạy bước tìm đường đi nào sau khi đã có kết quả;
mọi số trong lời giải thích được ghi lại trong cùng lần tính điểm.
\begin{itemize}
\item \textbf{VG, mô hình chính.} Vì trộn điểm thô, điểm cuối tách chính xác thành
$a\cdot\cos + \sum_i b\,w_i \sum_j a_{ij}\mathit{sim}_{ij} + \sum_i c\,w_i \sum_j a_{ij}\mathit{sim}_{ij}$. Bảng giải thích
ghi với mỗi phần $i$ của câu: $w_i$, phần $j$ của ảnh có attention lớn nhất, độ khớp, phần trăm của điểm cuối và, với vật
thể có số lượng, giá trị $g$ đã cộng. Một cột nữa ghi giá trị của dòng đó trừ trung bình của chính dòng đó trên 50 ứng
viên. Tổng cột này đúng bằng điểm cuối trừ điểm trung bình, nên cột này cho biết phần đóng góp của từng dòng vào chênh lệch giữa ảnh này và các ảnh
khác.
\item \textbf{MKG-W.} Lời giải thích là danh sách thực thể khớp được trong câu và các cạnh KGC-train giữa chúng với ứng viên.
Đây chính là các cạnh đã được đếm vào kênh \texttt{triple}.
\end{itemize}

% =====================================================================
\section{Kết quả cấp độ 1: Visual Genome $\cap$ COCO}
\label{sec:l1}

R@$k$ là tỉ lệ truy vấn có ảnh đúng trong $k$ kết quả đầu; MRR là trung bình nghịch đảo hạng của ảnh đúng, tính trong
top-50. Mọi số của mô hình huấn luyện là trung bình $\pm$ độ lệch chuẩn qua 3 seed, lấy ở checkpoint tốt nhất theo val.
Khi so hai mô hình, nhóm dùng kiểm định Welch trên 3 seed mỗi bên.

\subsection{So sánh CLIP với CLIP + đồ thị trên test}

\begin{table}[t]
\centering
\caption{Kết quả trên test VG $\cap$ COCO (10\,696 truy vấn, pool 2\,138 ảnh). Test chấm một lần cho mỗi mô hình đã chốt.}
\label{tab:l1-test}
\footnotesize
\setlength{\tabcolsep}{3pt}
\resizebox{\columnwidth}{!}{%
\begin{tabular}{lcccc}
\toprule
Phương pháp & R@1 & R@5 & R@10 & MRR \\
\midrule
CLIP thuần & 39,32 & 66,39 & 77,32 & 52,00 \\
Ba kênh, không huấn luyện & 42,54 & 70,04 & 80,35 & 55,05 \\
Mô hình cấp 1 (\texttt{main\_r3}) & \pmstd{45,77}{0,23} & \pmstd{72,87}{0,31} & \pmstd{82,69}{0,31} & \pmstd{58,14}{0,18} \\
\textbf{Mô hình chính} & \textbf{\pmstd{47,95}{0,15}} & \textbf{\pmstd{74,79}{0,03}} & \textbf{\pmstd{84,38}{0,08}} & \textbf{\pmstd{60,15}{0,15}} \\
\bottomrule
\end{tabular}}
\tabsrc{\file{experiments/level1/test/}, \file{experiments/improve/test/}}
\end{table}

Bảng~\ref{tab:l1-test} là phép so bắt buộc của đề: CLIP thuần và CLIP + đồ thị trên cùng cách chia. Dòng ``ba kênh, không
huấn luyện'' dùng cùng công thức trộn nhưng vector của phần đồ thị là vector CLIP text gốc, không qua GAT. Mô hình chính hơn
CLIP 8,63 điểm R@1. So với mô hình cấp~1, mô hình chính hơn 2,18 điểm R@1 ($p = 0{,}0005$) và 2,00 điểm MRR
($p = 0{,}0002$). Ba seed của mô hình chính cho R@1 48,06 / 47,78 / 47,99. Từ val sang test, R@1 giảm 0,96 điểm (48,91
$\to$ 47,95), so với 0,71 điểm của mô hình cấp~1; mô hình chính được xây qua nhiều bước chọn trên val
(Bảng~\ref{tab:l1-chain}).

\subsection{Mô hình chính được xây thế nào}

\begin{table}[t]
\centering
\caption{Từng bước từ mô hình cấp~1 tới mô hình chính, R@1 và MRR trên val. $p$: Welch so với bước trước.}
\label{tab:l1-chain}
\footnotesize
\setlength{\tabcolsep}{2.5pt}
\begin{tabular}{p{3.4cm}ccc}
\toprule
Bước (đổi một thứ so với bước trước) & R@1 & MRR & Chênh R@1 ($p$) \\
\midrule
Mô hình cấp 1 & \pmstd{46,48}{0,24} & 58,85 & \\
+ Đồ thị câu FACTUAL-T5 & \pmstd{47,23}{0,36} & 59,49 & +0,76 ($\approx$0,05) \\
+ Trộn điểm thô & \pmstd{47,48}{0,10} & 59,72 & +0,25 (n.s.) \\
+ Đồ thị ảnh sạch và liền\textsuperscript{*} & \pmstd{48,38}{0,17} & 60,51 & +0,90 (0,003) \\
+ Số lượng và bảng $g$ & \textbf{\pmstd{48,91}{0,21}} & \textbf{60,86} & +0,53 (0,032) \\
\bottomrule
\multicolumn{4}{p{0.95\columnwidth}}{\scriptsize \textsuperscript{*}Gộp bốn thay đổi (gộp trùng, hình học cạnh, cạnh
kNN, node tổng) theo quyết định của nhóm; ablation bỏ cạnh ở Bảng~\ref{tab:l1-ablation-main} đo phần đóng góp của cạnh.}
\end{tabular}
\tabsrc{\file{experiments/improve/README.md}}
\end{table}

Mô hình chính được xây qua bốn bước; mỗi bước đổi một thứ và được giữ lại khi R@1 val tăng qua 3 seed
(Bảng~\ref{tab:l1-chain}). Mức tăng của bước trộn điểm thô không có ý nghĩa thống kê. Bước này vẫn được giữ vì nó cho phép
tách chính xác điểm cuối thành từng phần (Mục~\ref{sec:method-expl}) và có độ lệch giữa các seed nhỏ nhất. Các hướng thử trên val nhưng không
giữ (so khớp bộ ba tách thành phần, kênh mức cả đồ thị, gom vật cùng loại thành một node nhóm, \dots) ghi ở
\file{experiments/improve/README.md}. Ví dụ, gom vật cùng loại làm R@1 giảm 0,63 điểm ($p = 0{,}019$), nên mô hình chính chỉ
thêm số lượng vào từng node.

\subsection{Ablation: cải thiện có đến từ cấu trúc quan hệ không}

Đề yêu cầu bỏ cạnh hoặc thay bằng cạnh ngẫu nhiên để kiểm tra phần cải thiện có đến từ cấu trúc quan hệ. Nhóm làm
hai loại ablation. Loại thứ nhất huấn luyện lại từ đầu, cùng công thức, cùng 3 seed. Loại thứ hai giữ nguyên checkpoint và
chỉ phá đồ thị lúc chấm.

\begin{table}[t]
\centering
\caption{Ablation của mô hình cấp~1, mỗi dòng huấn luyện lại từ đầu. R@1.}
\label{tab:l1-ablation}
\footnotesize
\setlength{\tabcolsep}{3pt}
\begin{tabular}{p{4.3cm}cc}
\toprule
Cấu hình & Val & Test \\
\midrule
Mô hình cấp 1 & \pmstd{46,48}{0,24} & \pmstd{45,77}{0,23} \\
$a,b,c$ cố định cho mọi câu & \pmstd{46,13}{0,15} & \pmstd{45,13}{0,05} \\
Bỏ GAT (không lan truyền) & \pmstd{46,71}{0,07} & \pmstd{45,95}{0,20} \\
Tắt kênh bộ ba (GAT vẫn dùng cạnh) & \pmstd{45,99}{0,05} & \pmstd{45,07}{0,22} \\
Bỏ hẳn quan hệ & \pmstd{45,98}{0,10} & \pmstd{45,26}{0,14} \\
Câu không qua bộ mã hoá đồ thị & \pmstd{46,92}{0,07} & \pmstd{45,61}{0,29} \\
Nối lại cạnh ngẫu nhiên (giữ bậc) & \pmstd{46,24}{0,24} & \pmstd{45,50}{0,40} \\
Loss LambdaRank thay cross-entropy & \pmstd{45,18}{0,34} & \pmstd{44,28}{0,30} \\
\bottomrule
\end{tabular}
\tabsrc{\file{experiments/level1/README.md}}
\end{table}

\begin{table}[t]
\centering
\caption{Ablation và đối chứng của mô hình chính, chỉ trên val. R@1 (chênh so với mô hình đủ).}
\label{tab:l1-ablation-main}
\footnotesize
\setlength{\tabcolsep}{3pt}
\begin{tabular}{p{3.5cm}cc}
\toprule
 & Mô hình chính & Mô hình cấp 1 \\
\midrule
Đủ đồ thị & \pmstd{48,91}{0,21} & \pmstd{46,48}{0,24} \\
\multicolumn{3}{l}{\emph{Huấn luyện lại}} \\
Bỏ mọi cạnh & $-$2,39 ($p$ = 0,0024) & $-$0,50 ($p$ = 0,055) \\
\multicolumn{3}{l}{\emph{Chỉ phá lúc chấm}} \\
Đồ thị của một ảnh khác & 21,91 ($-$27,00) & 26,33 ($-$20,15) \\
Nối lại quan hệ có nhãn ngẫu nhiên & 48,32 ($-$0,59) & 46,29 ($-$0,19) \\
Bỏ bảng $g$ & $-$0,35 & --- \\
\bottomrule
\end{tabular}
\tabsrc{\file{experiments/improve/no-edges/REPORT.md}, \file{experiments/improve/controls/README.md}}
\end{table}

Bảng~\ref{tab:l1-ablation} và~\ref{tab:l1-ablation-main} cho ba kết luận.
\begin{itemize}
\item \textbf{Mô hình dựa vào đồ thị của đúng ảnh.} Khi mỗi ảnh mang scene graph của một ảnh khác, R@1 val của mô hình chính
giảm từ 48,91 xuống 21,91, thấp hơn CLIP thuần (39,13).
\item \textbf{Ở mô hình cấp~1, cạnh đóng góp ít.} Bỏ hẳn quan hệ làm R@1 giảm khoảng 0,5 điểm; bỏ GAT không làm giảm kết quả.
\item \textbf{Ở mô hình chính, cạnh đóng góp 2,39 điểm.} Huấn luyện lại sau khi bỏ mọi cạnh làm R@1 val giảm 2,39 điểm
($p = 0{,}0024$), tương ứng 24\% mức chênh giữa mô hình chính và CLIP trên val (2,39 / 9,78); ở mô hình cấp~1 tỉ lệ này là
7\%. Khi bỏ cạnh, mô hình chính còn 46,52, ngang mô hình cấp~1 (46,48). Xáo các quan hệ có nhãn lúc chấm làm R@1 giảm 0,59
điểm. Phần đóng góp riêng của từng loại cạnh (quan hệ có nhãn, kNN, node tổng) chưa được tách.
\end{itemize}

\subsection{Độ nhạy theo $\alpha$ và đồ thị giúp ở đâu}

\begin{figure}[t]
\centering
\begin{tikzpicture}
\begin{axis}[width=\columnwidth, height=4.6cm, xlabel={$\alpha$ (trọng số CLIP)}, ylabel={R@1 val (\%)},
  xmin=0, xmax=1, ymin=20, ymax=52, grid=major, grid style={gray!25},
  legend style={font=\scriptsize, at={(0.02,0.03)}, anchor=south west}, tick label style={font=\scriptsize},
  label style={font=\small}]
\addplot[cObj, thick, mark=*, mark size=1.2pt] coordinates {(0,32.81) (0.1,36.77) (0.2,40.73) (0.3,44.12) (0.4,46.52)
  (0.5,47.82) (0.6,48.14) (0.7,47.37) (0.8,45.49) (0.9,42.84) (1,39.13)};
\addlegendentry{VG $\cap$ COCO, mô hình chính}
\addplot[cTri, thick, mark=square*, mark size=1.1pt] coordinates {(0,32.40) (0.05,32.40) (0.1,32.40) (0.15,32.40)
  (0.2,32.40) (0.25,32.40) (0.3,32.40) (0.35,32.40) (0.4,32.06) (0.45,31.84) (0.5,31.61) (0.55,31.50) (0.6,30.61)
  (0.65,29.93) (0.7,29.71) (0.75,29.26) (0.8,28.36) (0.85,27.91) (0.9,26.57) (0.95,25.67) (1,24.33)};
\addlegendentry{MKG-W, CLIP + đếm cạnh ($\beta=0$)}
\end{axis}
\end{tikzpicture}
\caption{R@1 val khi ép một $\alpha$ cố định cho mọi truy vấn ($\alpha=1$ là CLIP thuần). Với VG, $b$ và $c$ giữ tỉ lệ đã
học của từng câu, trung bình 3 seed. Với MKG-W, seed 0 (đường cong giống hệt ở cả 3 seed vì $\beta=0$ không dùng NativE).
Nguồn: \file{experiments/improve/alpha-groups/alpha_sweep.json}, \file{experiments/level2/val/alpha_beta_sweep.json}.}
\label{fig:alpha}
\end{figure}

\textbf{Độ nhạy theo $\alpha$} (Hình~\ref{fig:alpha}). Trên VG, đường cong có đỉnh ở $\alpha = 0{,}6$ (R@1
\pmstd{48,14}{0,30}). Trọng số học theo từng câu cho 48,91, cao hơn $\alpha$ cố định tốt nhất 0,77 điểm, dù $\alpha$ cố định
này được chọn trên chính val. Trung bình $a$ học được là \pmstd{0,57}{0,06}, gần $\alpha$ cố định tốt nhất; mức chênh 0,77
điểm đến từ việc trọng số thay đổi theo từng câu. Chỉ dùng đồ thị ($\alpha = 0$) đạt 32,81 và chỉ dùng CLIP ($\alpha = 1$) đạt
39,13, đều thấp hơn mọi giá trị trong khoảng $\alpha \in [0{,}4;\,0{,}8]$. Đường cong của MKG-W được phân tích ở
Mục~\ref{sec:l2}.

\begin{table}[t]
\centering
\caption{R@1 val theo nhóm câu. Nhóm định nghĩa một lần từ đồ thị FACTUAL của câu, dùng chung cho mọi mô hình. KTC 95\%:
bootstrap ghép cặp theo câu.}
\label{tab:l1-groups}
\footnotesize
\setlength{\tabcolsep}{2pt}
\begin{tabular}{lrcccc}
\toprule
Nhóm & Số câu & CLIP & Cấp 1 & Chính & Chính $-$ cấp 1 [KTC] \\
\midrule
Tất cả & 10\,633 & 39,13 & 46,48 & 48,91 & +2,43 [1,92; 2,93] \\
Có số lượng nêu rõ & 1\,235 & 44,21 & 50,69 & 54,06 & +3,37 [1,75; 4,91] \\
Có $\ge$1 quan hệ & 10\,342 & 39,31 & 46,83 & 49,38 & +2,55 \\
Chỉ vật thể & 291 & 32,99 & 33,91 & 32,19 & $-$1,72 [$-$4,7; 1,4] \\
2 vật thể & 3\,757 & 35,67 & 40,87 & 42,41 & +1,54 [0,72; 2,40] \\
$\ge$3 vật thể & 6\,607 & 41,40 & 50,20 & 53,32 & +3,11 [2,45; 3,77] \\
\bottomrule
\end{tabular}
\tabsrc{\file{experiments/improve/alpha-groups/REPORT.md}}
\end{table}

\textbf{Đồ thị giúp ở đâu} (Bảng~\ref{tab:l1-groups}). Đây là câu hỏi của nhiệm vụ~2 trong đề.
\begin{itemize}
\item Mức tăng so với CLIP tăng theo số vật thể trong câu: +6,74 điểm ở câu có 2 vật thể, +11,92 điểm ở câu từ 3 vật thể
trở lên.
\item Với câu chỉ nêu vật thể, không có quan hệ (2,7\% số câu), mô hình chính thấp hơn CLIP 0,80 điểm.
\item Với câu nêu số lượng, mức tăng so với mô hình cấp~1 là +3,37 điểm; với các câu còn lại là +2,31 điểm.
\item Phần mất khi bỏ cạnh cũng tập trung ở câu có quan hệ và câu nhiều vật thể; câu chỉ có vật thể gần như không đổi
(\file{experiments/improve/no-edges/REPORT.md}).
\end{itemize}

\subsection{Giải thích và phân tích 10 truy vấn}
\label{sec:l1-cases}

\textbf{Cách chọn.} Chỉ xét câu val có ít nhất một quan hệ. Nhóm bốc ngẫu nhiên với seed cố định (2026), theo quy tắc đặt
trước: 5 câu thành công (CLIP xếp ảnh đúng dưới hạng 1, mô hình xếp hạng 1), 3 câu mô hình làm tệ đi (CLIP hạng 1, mô hình
dưới hạng 1), 2 câu cả hai đều sai. Trên toàn val, mô hình chính (seed 0) có 1\,528 câu thành công, 463 câu làm tệ đi,
4\,944 câu cả hai đều sai và 3\,698 câu cả hai đều đúng. Bảng~\ref{tab:l1-cases} tóm tắt 10 câu. Ảnh, bảng giải thích đầy
đủ và lời phân tích từng câu (viết sau khi xem ảnh thật) ở \file{experiments/improve/explanations/README.md}.

\begin{table*}[t]
\centering
\caption{Mười truy vấn của mô hình chính (seed 0, val). Hạng ảnh đúng theo CLIP $\to$ theo mô hình.}
\label{tab:l1-cases}
\footnotesize
\begin{tabular}{cp{7.6cm}cp{6.0cm}}
\toprule
Ca & Truy vấn & Hạng & Nhận xét \\
\midrule
1 & Dining table for six in a wood paneled room with a fireplace. & 2 $\to$ 1 & Thành công \\
2 & A cow standing along side of a large river. & 8 $\to$ 1 & Thành công \\
3 & A man sitting in front of the pen with a dog herding sheep in the background. & 3 $\to$ 1 & Thành công \\
4 & Someone holding a doughnut that has a bite taken out of it. & 4 $\to$ 1 & Thành công \\
5 & People huddle under several umbrellas on a sidewalk area and a white dog \dots & 2 $\to$ 1 & Thành công; $g$ dùng số lượng của \emph{umbrella} và \emph{people} \\
6 & A bunch of bananas are sitting on a wooden plate. & 1 $\to$ 2 & VG gọi đĩa gỗ là \emph{tray}; thuộc tính ``wooden'' chưa được dùng \\
7 & A train riding along the tracks along a bank of dirt. & 1 $\to$ 5 & Đồ thị VG của ảnh đúng rỗng (0 node) \\
8 & A bus that is reading ``89 Portola Plaza'' is traveling on. & 1 $\to$ 3 & Chi tiết quyết định là chữ trên biển; scene graph không biểu diễn được \\
9 & Two men standing in a living room playing video game with remotes. & 13 $\to$ 17 & Số lượng khớp đúng ($g$ +0,071 cho ``two men'') nhưng ảnh đúng thiếu nhãn tay cầm \\
10 & A batter is preparing to hit the ball at a game. & 17 $\to$ 28 & Ảnh đúng thiếu nhãn \emph{batter}, \emph{ball} \\
\bottomrule
\end{tabular}
\end{table*}

\textbf{Phép thử xoá.} Đây là mức fidelity cơ bản của cấp~1. Trên ảnh mô hình xếp hạng 1 của mỗi ca, nhóm xoá phần ảnh có
đóng góp lớn nhất vào điểm, rồi so với trung bình 10 lần xoá ngẫu nhiên một phần cùng loại.
\begin{itemize}
\item Ở 10/10 ca, xoá phần lớn nhất làm điểm giảm nhiều hơn xoá ngẫu nhiên.
\item Ở 7/10 ca, riêng việc xoá phần đó đã làm ảnh rơi khỏi hạng 1 (mô hình cấp~1: 5/10). Khi phần bị xoá là vật thể, 6/6 ca
rơi hạng; khi là bộ ba, chỉ 1/4.
\item Attention trung bình của cặp khớp nhất là 0,51 ở kênh vật thể và 0,17 ở kênh bộ ba, tức attention của kênh bộ ba
phân tán hơn.
\end{itemize}
Mục~\ref{sec:xai} đo fidelity, validity và sparsity trên toàn bộ val.

\textbf{Nguyên nhân thất bại} quan sát được khi xem ảnh thật, ở cả hai mô hình: chú thích VG thiếu nhãn hoặc đồ thị rỗng;
khác biệt từ vựng giữa caption COCO và nhãn VG (\emph{plate}/\emph{tray}, \emph{child}/\emph{boy}); thuộc tính (màu, chất
liệu) chưa được dùng; chi tiết mà scene graph không biểu diễn được (chữ trên biển hiệu); ảnh đúng nằm ngoài top-50 của CLIP.
Về thuộc tính, một phép thử không huấn luyện chỉ tăng R@1 val 0,36 điểm, dưới ngưỡng 0,5 điểm nhóm đặt trước, nên nhóm
không thêm thuộc tính vào mô hình (\file{experiments/checks/attribute-probe.REPORT.md}).

% =====================================================================
\section{Kết quả cấp độ 2: MKG-W và NativE}
\label{sec:l2}

\subsection{Bước a: tái lập link prediction của NativE}

\begin{table}[t]
\centering
\caption{Link prediction trên tập test KGC của MKG-W (filtered). Bài báo chỉ công bố MRR và Hits@1 cho MKG-W. Dòng rmask:
3 seed.}
\label{tab:l2-repro}
\footnotesize
\setlength{\tabcolsep}{3pt}
\resizebox{\columnwidth}{!}{%
\begin{tabular}{lcccc}
\toprule
 & MRR & Hits@1 & Hits@3 & Hits@10 \\
\midrule
Bài báo NativE & 36,58 & 29,56 & --- & --- \\
Tái lập, MKG-W gốc (seed 42) & 36,99 & 30,17 & 39,95 & 49,32 \\
MKG-W-rmask (seed 0, 1, 2) & \pmstd{36,99}{0,24} & \pmstd{29,98}{0,11} & \pmstd{40,30}{0,33} & \pmstd{49,73}{0,54} \\
\bottomrule
\end{tabular}}
\tabsrc{\file{experiments/level2/native-repro/metrics.json}, \file{experiments/level2/native-train/}}
\end{table}

Nhóm chạy lại NativE trên MKG-W gốc với đúng mã nguồn và siêu tham số của tác giả. Tiêu chí đặt trước khi chạy: chênh không
quá khoảng 1 điểm thì coi là tái lập được. Kết quả lệch +0,41 MRR và +0,61 Hits@1 (Bảng~\ref{tab:l2-repro}), nên bước a đạt.
Bản \texttt{rmask}, đã che đặc trưng ảnh và chữ của 1\,784 thực thể val/test retrieval, cho cùng mức MRR, tức việc che
không làm giảm chất lượng link prediction. Mỗi lần huấn luyện mất khoảng 1,1 giờ trên một GPU RTX 4090.

Về độ lệch chuẩn: tham số \texttt{-seed} của mã nguồn chỉ đặt seed cho PyTorch; bộ lấy mẫu âm viết bằng C không được seed,
nên chuỗi mẫu âm giống nhau ở mọi seed. Độ lệch chuẩn qua 3 seed vì vậy chỉ phản ánh khởi tạo và nhiễu của phần đối kháng,
không bao gồm biến thiên do lấy mẫu âm.

\subsection{Bước b: kết quả truy vấn}

\begin{table}[t]
\centering
\caption{Kết quả truy vấn trên MKG-W (892 truy vấn mỗi split, pool 892 ảnh). $\alpha,\beta$ chọn trên val; test chấm một lần.
Ba dòng \texttt{kg\_triple}, \texttt{native}, \texttt{native\_kgc} trùng nhau vì val chọn $\beta = 0$.}
\label{tab:l2-results}
\footnotesize
\setlength{\tabcolsep}{2.5pt}
\resizebox{\columnwidth}{!}{%
\begin{tabular}{llcccc}
\toprule
 & Dòng & R@1 & R@5 & R@10 & MRR \\
\midrule
\multirow{5}{*}{\rotatebox{90}{Val}} & CLIP thuần & 24,33 & 44,39 & 53,81 & 34,22 \\
 & + đếm cạnh (dòng chính) & \textbf{32,40} & \textbf{50,78} & \textbf{57,62} & \textbf{41,18} \\
 & + điểm NativE, bỏ cạnh & \pmstd{28,25}{0,11} & \pmstd{45,85}{0,51} & \pmstd{54,37}{0,70} & \pmstd{36,97}{0,08} \\
 & + cosine NativE, bỏ cạnh & \pmstd{26,64}{0,26} & \pmstd{45,96}{0,39} & \pmstd{54,30}{0,62} & \pmstd{35,98}{0,15} \\
 & Cạnh xáo trộn & \pmstd{21,82}{0,39} & \pmstd{42,56}{0,06} & \pmstd{51,72}{0,26} & \pmstd{32,03}{0,23} \\
\midrule
\multirow{5}{*}{\rotatebox{90}{Test}} & CLIP thuần & 24,89 & 44,73 & 54,48 & 34,61 \\
 & + đếm cạnh (dòng chính) & \textbf{29,82} & \textbf{51,46} & \textbf{59,42} & \textbf{39,89} \\
 & + điểm NativE, bỏ cạnh & \pmstd{27,24}{0,19} & \pmstd{46,52}{0,11} & \pmstd{56,28}{0,11} & \pmstd{36,67}{0,14} \\
 & + cosine NativE, bỏ cạnh & \pmstd{26,08}{0,45} & \pmstd{44,92}{0,53} & \pmstd{54,22}{0,95} & \pmstd{35,46}{0,58} \\
 & Cạnh xáo trộn & \pmstd{21,45}{0,06} & \pmstd{43,31}{0,51} & \pmstd{53,29}{0,13} & \pmstd{32,02}{0,16} \\
\bottomrule
\end{tabular}}
\tabsrc{\file{experiments/level2/val/summary.json} và \file{experiments/level2/test/summary.json}}
\end{table}

Bảng~\ref{tab:l2-results} gom các dòng theo vai trò. ``Đếm cạnh'' là kênh \texttt{triple}; ``điểm NativE'' là kênh
\texttt{kgc} (dòng \texttt{kgc\_only}, $\beta = 1$); ``cosine NativE'' là kênh \texttt{emb} (dòng \texttt{native\_emb},
$\beta = 1$). Dòng ``cạnh xáo trộn'' đổi chỗ đuôi của các triple trong cùng một quan hệ, giữ nguyên bậc của đầu, rồi dùng
lại $\alpha,\beta$ của dòng chính. Dòng chính hơn CLIP 4,93 điểm R@1 và 5,28 điểm MRR trên test.

\textbf{Thứ tự kênh kgc được thêm.} Lần chạy val đầu tiên chỉ có kênh cosine, và val chọn $\beta = 0$ cho kênh này. Vì kênh
cosine không dùng quan hệ, kênh \texttt{kgc} được thêm để dùng hàm chấm có quan hệ của chính NativE. Định
nghĩa kênh và dòng chính mới được chốt trước khi chạy lại val và trước lần chấm test ngày 17/11. Lần chạy lại val giữ
nguyên từng byte các file kết quả cũ (\file{experiments/level2/val/README.md}).

\textbf{Các lần đọc test khác và một lần chạy lại độc lập.} (1) Khi kiểm tra code đánh giá,
một lần chạy thử \texttt{--split test --hand} với embedding giả đã in ra số test của dòng đếm cạnh, kể cả trên bộ câu tự
viết. Dòng này không có tham số nào ngoài $\alpha$, và $\alpha$ đã chốt trên val; không có gì được chỉnh sau đó.
(2) Ngày 18/11, Khang chạy lại độc lập toàn bộ cấp~2 trên molab: tự huấn luyện NativE 3 seed (MRR rmask \pmstd{36,89}{0,27},
bản gốc 37,06), dùng đặc trưng CLIP gốc thay cho bản nhóm trích lại, tự chọn $\alpha,\beta$ trên val rồi chấm test một lần.
Val lại chọn $\alpha = 0$, $\beta = 0$ ở cả 3 seed; dòng chính trên test cho đúng R@1 29,82, cosine NativE
\pmstd{26,05}{0,51} và cạnh xáo trộn \pmstd{21,41}{0,11}, gần với Bảng~\ref{tab:l2-results}. CLIP thuần trên test là 25,00
thay vì 24,89 vì khác bản đặc trưng. Kết quả và log của lần chạy này ở \file{experiments/level2/molab-rerun/}.

\subsection{Phân tích kết quả}

\begin{itemize}
\item \textbf{Kênh đếm cạnh cho mức tăng lớn nhất.} Kênh này không có tham số huấn luyện và nâng R@1 val từ 24,33 lên 32,40.
\item \textbf{Đối chứng xáo cạnh.} Xáo đuôi cạnh làm R@1 thấp hơn CLIP (val 21,82, test 21,45). Vì phép xáo giữ nguyên bậc
của đầu cạnh, mức tăng của dòng chính đến từ việc cạnh nối đúng thực thể, không đến từ bậc của thực thể.
\item \textbf{NativE dùng riêng.} Điểm bộ ba của NativE nâng R@1 val 3,92 điểm so với CLIP; cosine embedding nâng 2,32 điểm.
Thứ tự này giữ nguyên trên test và ở cả 3 seed: kênh dùng quan hệ (\texttt{kgc}) cho kết quả cao hơn kênh chỉ so vector thực
thể (\texttt{emb}), nhưng cả hai thấp hơn kênh đếm cạnh.
\item \textbf{NativE kết hợp với đếm cạnh.} Ở cả 3 seed, val chọn $\beta = 0$ cho cả kênh cosine lẫn kênh \texttt{kgc}. Mọi
$\beta > 0$ đều cho R@1 val thấp hơn (tốt nhất 31,61 với cosine và 30,38 với \texttt{kgc}, so với 32,40). Do kênh NativE có
trọng số 0, thay embedding của nó bằng nhiễu Gauss cho đúng con số của dòng chính. Như vậy, trong dòng chính của cấp~2,
NativE không đóng góp; kết quả của NativE được báo cáo ở phần tái lập (Bảng~\ref{tab:l2-repro}) và ở hai dòng dùng riêng.
\end{itemize}

\textbf{Độ nhạy theo $\alpha$} (Hình~\ref{fig:alpha}). Trên MKG-W, R@1 val giữ ở 32,40 trong cả đoạn $\alpha \in [0; 0{,}35]$
rồi giảm dần về 24,33 ở $\alpha = 1$; giá trị được chọn $\alpha = 0$ nằm trong vùng phẳng này. Khác với VG (đỉnh ở
$\alpha = 0{,}6$), trên MKG-W chỉ dùng đồ thị (32,40) cho kết quả cao hơn chỉ dùng CLIP (24,33). Với $\alpha = 0$, điểm cuối
chỉ còn số cạnh; các ứng viên cùng số cạnh giữ thứ tự của CLIP vì phép sắp xếp ổn định.

\textbf{Theo nhóm truy vấn.} Truy vấn không khớp được thực thể nào (189 câu val, 170 câu test) giữ nguyên kết quả của CLIP
theo thiết kế. Toàn bộ mức tăng nằm ở nhóm khớp được: R@1 val 20,91 $\to$ 31,15, test 22,71 $\to$ 28,81.

\subsection{Bộ 25 truy vấn tự viết: trực tiếp và gián tiếp}
\label{sec:l2-hand}

\textbf{Cách xây và kiểm tra.} Bộ câu nằm ở \file{data/queries/mkgw_queries_hand.json}. Quy tắc được kiểm tra tự động bởi
\file{scripts/build_mkgw_hand_queries.py}; script từ chối ghi file nếu có câu vi phạm.
\begin{itemize}
\item \textbf{Trực tiếp} (10 câu): tả thực thể đáp án là gì, trông ra sao, không nêu tên nó và không dựa vào thực thể có tên
khác. Mỗi câu có đúng một đáp án trong kho ảnh test. Script kiểm tra mọi thực thể khớp được trong câu đều không có cạnh tới
đáp án, nên câu trực tiếp không có đường tắt qua đồ thị.
\item \textbf{Gián tiếp} (15 câu): nêu tên một thực thể khác và cần một quan hệ để tới đáp án, ví dụ ``a rock guitarist who
was a member of Pink Floyd''. Đáp án không gõ tay mà được tính từ mẫu (quan hệ, thực thể) của câu: mọi thực thể trong kho
ảnh test có đủ các cạnh đó, từ 1 tới 5 đáp án.
\item Tên đáp án không được xuất hiện trong câu. Đáp án tính trên toàn bộ triple MKG-W, vì đó là đáp án thật chứ không phải
đầu vào của mô hình. Kênh đồ thị lúc chấm vẫn chỉ thấy cạnh KGC-train. Bộ câu được commit một lần (16/11) và không sửa sau đó.
\end{itemize}

\begin{table}[t]
\centering
\caption{Bộ 25 truy vấn tự viết, ảnh của tập test làm ứng viên, dùng lại $\alpha,\beta$ của val.}
\label{tab:l2-hand}
\footnotesize
\setlength{\tabcolsep}{3pt}
\begin{tabular}{llcccc}
\toprule
Nhóm & Dòng & R@1 & R@5 & R@10 & MRR \\
\midrule
Trực tiếp (10) & CLIP & 80,00 & 90,00 & 90,00 & 85,00 \\
 & + đếm cạnh & 80,00 & 90,00 & 90,00 & 85,00 \\
Gián tiếp (15) & CLIP & 40,00 & 60,00 & 60,00 & 49,28 \\
 & + đếm cạnh & \textbf{80,00} & \textbf{86,67} & \textbf{86,67} & \textbf{83,33} \\
\bottomrule
\end{tabular}
\tabsrc{\file{experiments/level2/test/summary.json} (khoá \texttt{hand\_*})}
\end{table}

\textbf{Kết quả} (Bảng~\ref{tab:l2-hand}). Với câu gián tiếp, kênh đếm cạnh nâng R@1 từ 40 lên 80 (6/15 $\to$ 12/15). Câu
trực tiếp không đổi: CLIP đã đúng 8/10, và theo cách xây, kênh đồ thị không cộng điểm cho ảnh đúng của các câu này. Nhóm tra
xem cạnh vàng của 15 câu gián tiếp nằm ở phần nào của KGC. 12/15 câu có ít nhất một đáp án nối với thực thể trong câu bằng
cạnh KGC-train, nên kênh đếm cạnh thấy được. 3/15 câu chỉ có cạnh vàng trong KGC-test, nên kênh này không thể giúp; ví dụ
``the rock band whose lead singer is Hayley Williams'' (đáp án Paramore). Do bộ câu chỉ có 25 câu, kết quả này mang tính
minh hoạ và không thay cho kết quả trên 892 truy vấn.

\textbf{Ví dụ lời giải thích.} Với câu ``the main railway station of Bordeaux'', \texttt{EntityLinker} khớp
\emph{Bordeaux}. Ảnh đúng là ga Bordeaux-Saint-Jean, có cạnh KGC-train P131 (\emph{located in the administrative territorial
entity}) tới Bordeaux. Lời giải thích của ứng viên là cạnh này, cũng là cạnh được đếm vào điểm đồ thị. Với câu về Pink
Floyd, có hai cạnh KGC-train được đếm: P463 (\emph{member of}, David Gilmour $\to$ Pink Floyd) và P527 (\emph{has part},
Pink Floyd $\to$ David Gilmour), nên điểm đồ thị của ảnh đúng là 2. \todo{Hạng của ảnh đúng theo CLIP và theo
dòng chính cho từng câu: chờ bảng theo câu (việc 8, Khoa); lần chạy 17/11 chỉ lưu số tổng theo nhóm.}

\textbf{Về GAT trên MKG-W.} MKG-W không có scene graph cho từng ảnh, nên bộ xếp lại GATv2 của cấp~1 không áp dụng được. Ở
bước b, tín hiệu từ đồ thị tri thức (đếm cạnh và NativE) thay vai trò của GNN. Trên MKG-W, phép so ``CLIP với CLIP + đồ thị''
của đề là hai dòng đầu của Bảng~\ref{tab:l2-results}.

\InputIfFileExists{sections/06-explanation-eval}{}{\section{Đánh giá lời giải thích}\label{sec:xai}\todo{chưa có mục}}
\InputIfFileExists{sections/07-vietnamese-mmkg}{}{\section{MMKG tiếng Việt}\label{sec:vnmmkg}\todo{chưa có mục}}
% =====================================================================
\section{Demo, hạn chế và kết luận}
\label{sec:conclusion}

\subsection{Demo}
\label{sec:demo}

\file{app/demo_app.py} là một giao diện Gradio, chạy được trên GPU 8\,GB. Người dùng nhập truy vấn tiếng Anh, chọn CLIP thuần
hoặc CLIP + đồ thị, rồi xem top-$k$ ảnh. Mỗi ảnh kèm bảng giải thích như Mục~\ref{sec:method-expl}, khung của các vật thể khớp
vẽ trên ảnh, và cột ``bỏ phần này'': chấm lại ảnh bằng chính mô hình sau khi bỏ phần khớp đó khỏi đồ thị ảnh. Câu gõ vào được
tách bằng FACTUAL-T5 và cùng luật số lượng như lúc đóng gói dữ liệu. Phép kiểm tra ở Mục~\ref{sec:checks} bảo đảm demo cho
cùng kết quả với pipeline đánh giá. Video dự phòng theo kịch bản \file{docs/demo-video-script.md}.

\subsection{Hạn chế}

\begin{itemize}
\item \textbf{Chỉ 3 seed.} Mọi kiểm định giữa các cấu hình chỉ dựa trên 3 seed mỗi bên.
\item \textbf{Val được dùng nhiều lần} để chọn từng bước của mô hình chính VG, nên số val hơi lạc quan (từ val sang test giảm
0,96 điểm R@1). Test vẫn chỉ chấm một lần.
\item \textbf{Chưa tách được đóng góp của từng loại cạnh} trong mô hình chính VG: bước đồ thị liền gộp bốn thay đổi, và chưa có
ablation chỉ bỏ cạnh kNN.
\item \textbf{Đếm số lượng coi mọi người là một loại.} Câu ``a lone woman'' trong ảnh nhiều người bị coi là sai số lượng, và
bảng $g$ đôi khi kéo phần khớp sang một node khác nhãn.
\item \textbf{Lỗi parser.} Thư viện FactualSceneGraph 0.7.3 cắt token cuối của câu dài nhất trong mỗi batch. Khoảng 1,8\%
caption bị ảnh hưởng; 0,26\% (train), 0,24\% (val) và 0,31\% (test) mất một phần chữ. Lỗi xảy ra như nhau ở cả ba split.
Nhóm không parse lại dữ liệu, giữ nguyên các số đã báo, và đã sửa lỗi này trong demo.
\item \textbf{Thuộc tính và chữ trong ảnh chưa được dùng}, dù nhiều ca thất bại liên quan tới chúng.
\item \textbf{Giới hạn của xếp lại.} 3,3\% câu test VG có ảnh đúng nằm ngoài top-50 của CLIP và không cứu được bằng xếp lại.
\item \textbf{MKG-W.} Phương pháp năm 2024 không đóng góp vào dòng chính ($\beta = 0$). Đặc trưng ảnh của NativE (383 chiều,
do tác giả trích sẵn) khác CLIP ViT-B/32 của nhóm. Cache Wikidata phải lấy lại ngày 17/11: 441 thực thể mất nhãn tiếng Anh,
số truy vấn val khớp được thực thể đổi từ 704 xuống 703. Bộ câu tự viết chỉ có 25 câu.
\end{itemize}
Hạn chế riêng của phần đánh giá lời giải thích và phần MMKG tiếng Việt ghi ở Mục~\ref{sec:xai} và Mục~\ref{sec:vnmmkg}.

\subsection{Kết luận}

Báo cáo trả lời ba câu hỏi của đề như sau.
\begin{enumerate}
\item \textbf{Kết quả nào phù hợp nhất?} Trên VG $\cap$ COCO, xếp lại top-50 của CLIP bằng scene graph nâng R@1 test từ
39,32 lên 47,95. Trên MKG-W, đếm cạnh tri thức nâng R@1 test từ 24,89 lên 29,82.
\item \textbf{Tri thức quan hệ giúp ở đâu?} Ở câu có nhiều vật thể và có quan hệ: phần lời tăng theo số vật thể, còn câu chỉ
có một vật thể thì không được lợi. Ở câu gián tiếp: trên 15 câu tự viết, R@1 tăng từ 40 lên 80, trong khi câu trực tiếp
không đổi. Các phép đối chứng cho thấy phần lời đến từ đúng cấu trúc: thay đồ thị của ảnh khác hoặc xáo cạnh tri thức đều
làm kết quả tụt dưới CLIP.
\item \textbf{Vì sao mô hình trả về kết quả này?} Lời giải thích là chính các đại lượng của điểm đồ thị. Với mô hình chính
của VG, điểm tách chính xác thành từng phần. Mục~\ref{sec:xai} đo mức trung thực của lời giải thích đó.
\end{enumerate}
Nhóm cũng ghi lại các kết quả âm có bằng chứng: ở cấp~1 cạnh đóng góp ít, và NativE không thêm được gì khi đã có kênh đếm
cạnh.

% =====================================================================
\section{Phân công và khai báo}
\label{sec:contrib}

\subsection{Bảng phân công}

Bảng~\ref{tab:contrib} tóm tắt phần việc của từng thành viên, lấy từ bảng nhận việc \file{task-assignment.xlsx} và lịch sử
commit. Minh chứng chi tiết theo từng dòng việc nằm trong file bảng nhận việc và lịch sử git của repo.
\todo{Cả nhóm xác nhận bảng này trước khi nộp; bảng nhận việc cần cập nhật cột Trạng thái.}

\begin{table*}[t]
\centering
\caption{Phân công đóng góp.}
\label{tab:contrib}
\footnotesize
\begin{tabular}{p{3.6cm}p{13.2cm}}
\toprule
Thành viên & Phần việc \\
\midrule
Nguyễn Đình Lộc (25C11050) &
Khung repo, trích đặc trưng CLIP, module đánh giá, chia tập VG $\cap$ COCO, baseline CLIP, kiểm tra pipeline; toàn bộ cấp
độ~1 (dựng đồ thị, mô hình GATv2, huấn luyện, ablation, giải thích, 10 truy vấn) và các bước xây mô hình chính; demo. Cấp
độ~2: huấn luyện NativE trên GPU 4090, xuất embedding, kênh \texttt{kgc}, dựng lại đặc trưng MKG-W, chạy đánh giá val/test.
Bản đo giải thích v1. Điều phối, review và merge; ráp báo cáo. \\
Ngô Trương Minh Đạt (25C15003) &
Dữ liệu MKG-W: lấy Wikidata, chia tập retrieval, sinh truy vấn che tên; \texttt{EntityLinker} và hàm điểm đồ thị của bước b;
hỗ trợ tái lập NativE. \\
Trần Đắc Khoa (25C15015) &
Phân tích kết quả và bản nháp báo cáo cấp độ~2 (tái lập, bảng val/test, ablation, độ nhạy $\alpha$, trực tiếp/gián tiếp);
ví dụ giải thích trên MKG-W. \\
Nguyễn Bùi Mẫn Nhi (25C15055) &
Cấp độ~3: định nghĩa và cài đặt fidelity, validity, sparsity (bản v2, sửa các lỗi của v1), so sánh các cách giải thích
(v3: phản thực tế, GraphMask, GStarX; ba benchmark có đáp án biết trước),
bằng chứng trên từng truy vấn, cột ``bỏ phần này'' của demo; mô hình scene ranker song song (Sinkhorn trên concept graph)
và báo cáo của nó. \\
Âu Dương Khang (25C15045) &
Cấp độ~3: MMKG tiếng Việt (chọn miền, thu thập, liên kết ảnh--thực thể, duyệt triple, truy vấn tiếng Việt, data card),
chạy chuyển giao; hỗ trợ chọn encoder đa ngữ. Cấp độ~2: bộ 25 truy vấn tự viết, script đánh giá MKG-W, notebook molab;
huấn luyện NativE 3 seed và chạy lại độc lập toàn bộ cấp~2 trên molab. \\
\bottomrule
\end{tabular}
\end{table*}

\subsection{Khai báo mã nguồn mở, dữ liệu và công cụ AI}

\textbf{Mã nguồn mở.} Nhóm dùng các thư viện qua API, không sửa mã nguồn của thư viện: PyTorch, PyTorch Geometric
(\texttt{GATv2Conv})~\cite{fey2019pyg}, open\_clip, FAISS, spaCy, SceneGraphParser, FactualSceneGraph, transformers, NLTK
(WordNet~\cite{miller1995wordnet}), Gradio. Mô hình xếp lại của cấp độ~1 do nhóm viết theo thiết kế riêng, không dùng mã
nguồn của bài báo nào. Ở cấp độ~2, nhóm chạy mã nguồn công khai của NativE (\texttt{zjukg/NativE}, commit \texttt{c72291c})
và chỉ sửa cách build: thư viện \texttt{Base.so} kèm sẵn bị lỗi segfault nên được build lại từ mã nguồn. Kênh \texttt{kgc}
tính lại hàm điểm của NativE trong code của nhóm và có test đối chiếu với bản gốc. Một lỗi của FactualSceneGraph được vá
trong code của nhóm, không sửa thư viện.

\textbf{Dữ liệu.} Visual Genome 1.4 (CC BY 4.0), caption MS-COCO 2014 (CC BY 4.0), split Karpathy, MKG-W và bộ ảnh của nhóm
tác giả MMRNS, Wikidata (CC0). Nguồn và giấy phép của MMKG tiếng Việt ghi ở Mục~\ref{sec:vnmmkg}. Link tải và checksum:
\file{data/README.md}, \file{data/MANIFEST.md}.

\textbf{Công cụ AI.} Nhóm dùng Anthropic Claude (Claude Code) và OpenAI Codex để lên phương án cùng người làm, viết spec,
viết phần lớn code theo spec đã duyệt, review code, chạy thực nghiệm và soạn nháp tài liệu. Người làm của từng phần chốt
thiết kế và mọi lựa chọn (trong đó có chọn mô hình chính trước khi chấm test), duyệt code và kết quả, và chịu trách nhiệm
giải thích phần của mình. Mọi thay đổi nằm trong lịch sử git, phần lớn qua pull request có người duyệt; số liệu chỉ lấy từ
file kết quả có log và seed. Bảng chi tiết phần nào do công cụ nào làm, theo từng phần việc: \file{report/declaration-ai-oss.md}.
\todo{Đạt, Khoa, Nhi bổ sung phần khai báo của mình vào file này (cấp độ~2, đo giải thích v2, scene ranker).}


\bibliographystyle{IEEEtran}
\bibliography{references,refs-xai,refs-vnmmkg}

\end{document}
